% ---------------------------------------------------------------- 2 related work
\section{Related Work}
\label{sec:related}

\subsection{Learned Semantic Transmission of Text}
\label{sec:related-jscc}

Deep joint source--channel coding (JSCC) replaces the separate source and channel codes by one learned mapping from content to channel symbols. For text, Farsad et al.\ \cite{farsad2018} encoded sentences with recurrent networks, and DeepSC \cite{xie2021deepsc} used a Transformer \cite{vaswani2017} with a fixed number of symbols per word, trained and evaluated on the Europarl corpus \cite{koehn2005}; its restriction to sentences of 4--30 words has been adopted by many later studies. For images, deep JSCC avoids the cliff effect of separate coding \cite{bourtsoulatze2019} and can deliver progressively refinable bandwidth-agile codes \cite{kurka2021}. Surveys and tutorials place these designs within the broader program of semantic and goal-oriented communication \cite{strinati2021,gunduz2023,luo2022,kountouris2021}. Two properties of this line of work matter here. The encoder and decoder are typically trained jointly for a data distribution that is assumed fixed, and the number of channel uses per message is either fixed or chosen from a small set at the transmitter.

Rate adaptation at the receiver's request has been studied as semantic hybrid automatic repeat request (HARQ). Jiang et al.\ \cite{jiang2022harq} combined a learned sentence codec with retransmission and incremental knowledge, following the incremental-redundancy principle of classical HARQ \cite{caire2001}. Successive refinement \cite{equitz1991} and nested representations \cite{kusupati2022} supply the structural idea behind prefix codes whose leading part is useful on its own. DN-SC applies incremental redundancy per token rather than per message and makes the request depend on the posterior confidence of the receiver's language prior; it also charges the requests on the return link, which a token-level policy makes frequent.

\subsection{Language Models in Semantic Communication}
\label{sec:related-llm}

Pre-trained language models enter semantic communication in three roles. As a measure of importance, they rank words so that power or bandwidth can be allocated to those that matter most \cite{guo2023importance}. As a knowledge base, large AI models supply the background knowledge of the transceiver and are argued to be a natural foundation for 6G semantic systems \cite{jiang2024largeai,xie2025largemodel,bariah2024}; LLM-based text transceivers have been proposed that dispense with a fixed knowledge base \cite{chen2025llmsc}, and large models have been fine-tuned and compressed for multi-task semantic communication \cite{ni2025multitask}. As a prior, finally, the language model provides the probability of the next token given what has been decoded, which the receiver combines with the channel likelihood. Wang et al.\ \cite{wang2024llmsc} built a text semantic communication system on this principle, with tokens as semantic units and a pre-trained LLM as the prior of the decoder; Qiao et al.\ \cite{qiao2025token} exploited the cross-modal context among tokens, through large models, to save bandwidth. DN-SC belongs to this third group and takes the LLM-prior decoder as its starting point. It adds the nested codebook with exact per-shell power control, token-level selective refinement, and memories that change after deployment.

\subsection{Knowledge Mismatch and Adaptation After Deployment}
\label{sec:related-adapt}

A semantic receiver works because it shares knowledge with the transmitter, and that knowledge can diverge. Zhang et al.\ \cite{zhang2023dynamic} studied transceivers that must cope with a task-unaware transmitter and dynamic data. Zhang et al.\ \cite{zhang2026kbident} test whether the knowledge bases of the two ends match and adapt them with a lightweight step when they do not, and cooperative update policies keep the knowledge bases of several semantic communication pairs current while limiting the overhead of the update \cite{zhang2026kbupdate}. In machine learning, adaptation to a changing data stream raises the problem of catastrophic forgetting \cite{kirkpatrick2017}; low-rank adaptation (LoRA) \cite{hu2022lora} makes updates of a large model small; and communication-efficient distributed learning compresses updates by quantization \cite{alistarh2017,mcmahan2017}. The nested-learning paradigm \cite{behrouz2025} organizes a model as components updated at different frequencies. In a transceiver, however, the receiver cannot compute its own updates without the source text, so every update has a price in air time that grows with its size and frequency. DN-SC differs from the work above in that it charges every update of a language-model prior, at three timescales, bit by bit and in channel symbols, against the rate that the update saves; the result is the traffic volume above which updating pays.

\subsection{Positioning}
\label{sec:related-position}

Table~\ref{tab:related} positions DN-SC against the closest work along the dimensions this article studies. The table records what each work addresses as its subject; a dash does not mean that the work could not be extended in that direction.

\begin{table}[t]
\centering
\caption{Positioning of DN-SC against representative work. \checkmark: addressed as a subject of the work; ---: not addressed as a subject.}
\label{tab:related}
\footnotesize
\begin{tabular}{@{}lccccc@{}}
\toprule
Work & \makecell{Pre-trained\\language model\\in the link} & \makecell{Receiver-driven\\incremental\\redundancy} & \makecell{Knowledge or\\model adaptation\\after deployment} & \makecell{Update overhead\\as a design\\quantity} & \makecell{Several\\update\\timescales} \\
\midrule
Semantic HARQ \cite{jiang2022harq}                 & ---        & \checkmark & ---        & ---        & --- \\
Importance-aware SC \cite{guo2023importance}       & \checkmark & ---        & ---        & ---        & --- \\
LLM-enabled text SC \cite{wang2024llmsc}           & \checkmark & ---        & ---        & ---        & --- \\
Token communications \cite{qiao2025token}          & \checkmark & ---        & ---        & ---        & --- \\
Dynamic-data SC \cite{zhang2023dynamic}            & ---        & ---        & \checkmark & ---        & --- \\
Knowledge-base adaptation \cite{zhang2026kbident}  & ---        & ---        & \checkmark & ---        & --- \\
Knowledge-base update \cite{zhang2026kbupdate}     & ---        & ---        & \checkmark & \checkmark & --- \\
\textbf{DN-SC (this work)}                         & \checkmark & \checkmark & \checkmark & \checkmark & \checkmark \\
\bottomrule
\end{tabular}
\end{table}
